\IfFileExists{stacks-project.cls}{%
\documentclass{stacks-project}
}{%
\documentclass{amsart}
}

% For dealing with references we use the comment environment
\usepackage{verbatim}
\newenvironment{reference}{\comment}{\endcomment}
%\newenvironment{reference}{}{}
\newenvironment{slogan}{\comment}{\endcomment}
\newenvironment{history}{\comment}{\endcomment}

% For commutative diagrams we use Xy-pic
\usepackage[all]{xy}

% We use 2cell for 2-commutative diagrams.
\xyoption{2cell}
\UseAllTwocells

% We use multicol for the list of chapters between chapters
\usepackage{multicol}

% This is generally recommended for better output
\usepackage{lmodern}
\usepackage[T1]{fontenc}

% For cross-file-references
\usepackage{xr-hyper}

% Package for hypertext links:
\usepackage{hyperref}

% For any local file, say "hello.tex" you want to link to please
% use \externaldocument[hello-]{hello}
\externaldocument[introduction-]{introduction}
\externaldocument[conventions-]{conventions}
\externaldocument[sets-]{sets}
\externaldocument[categories-]{categories}
\externaldocument[topology-]{topology}
\externaldocument[sheaves-]{sheaves}
\externaldocument[sites-]{sites}
\externaldocument[stacks-]{stacks}
\externaldocument[fields-]{fields}
\externaldocument[algebra-]{algebra}
\externaldocument[brauer-]{brauer}
\externaldocument[homology-]{homology}
\externaldocument[derived-]{derived}
\externaldocument[simplicial-]{simplicial}
\externaldocument[more-algebra-]{more-algebra}
\externaldocument[smoothing-]{smoothing}
\externaldocument[modules-]{modules}
\externaldocument[sites-modules-]{sites-modules}
\externaldocument[injectives-]{injectives}
\externaldocument[cohomology-]{cohomology}
\externaldocument[sites-cohomology-]{sites-cohomology}
\externaldocument[dga-]{dga}
\externaldocument[dpa-]{dpa}
\externaldocument[sdga-]{sdga}
\externaldocument[hypercovering-]{hypercovering}
\externaldocument[schemes-]{schemes}
\externaldocument[constructions-]{constructions}
\externaldocument[properties-]{properties}
\externaldocument[morphisms-]{morphisms}
\externaldocument[coherent-]{coherent}
\externaldocument[divisors-]{divisors}
\externaldocument[limits-]{limits}
\externaldocument[varieties-]{varieties}
\externaldocument[topologies-]{topologies}
\externaldocument[descent-]{descent}
\externaldocument[perfect-]{perfect}
\externaldocument[more-morphisms-]{more-morphisms}
\externaldocument[flat-]{flat}
\externaldocument[groupoids-]{groupoids}
\externaldocument[more-groupoids-]{more-groupoids}
\externaldocument[etale-]{etale}
\externaldocument[chow-]{chow}
\externaldocument[intersection-]{intersection}
\externaldocument[pic-]{pic}
\externaldocument[weil-]{weil}
\externaldocument[adequate-]{adequate}
\externaldocument[dualizing-]{dualizing}
\externaldocument[duality-]{duality}
\externaldocument[discriminant-]{discriminant}
\externaldocument[derham-]{derham}
\externaldocument[local-cohomology-]{local-cohomology}
\externaldocument[algebraization-]{algebraization}
\externaldocument[curves-]{curves}
\externaldocument[resolve-]{resolve}
\externaldocument[models-]{models}
\externaldocument[functors-]{functors}
\externaldocument[equiv-]{equiv}
\externaldocument[pione-]{pione}
\externaldocument[etale-cohomology-]{etale-cohomology}
\externaldocument[proetale-]{proetale}
\externaldocument[relative-cycles-]{relative-cycles}
\externaldocument[more-etale-]{more-etale}
\externaldocument[trace-]{trace}
\externaldocument[crystalline-]{crystalline}
\externaldocument[spaces-]{spaces}
\externaldocument[spaces-properties-]{spaces-properties}
\externaldocument[spaces-morphisms-]{spaces-morphisms}
\externaldocument[decent-spaces-]{decent-spaces}
\externaldocument[spaces-cohomology-]{spaces-cohomology}
\externaldocument[spaces-limits-]{spaces-limits}
\externaldocument[spaces-divisors-]{spaces-divisors}
\externaldocument[spaces-over-fields-]{spaces-over-fields}
\externaldocument[spaces-topologies-]{spaces-topologies}
\externaldocument[spaces-descent-]{spaces-descent}
\externaldocument[spaces-perfect-]{spaces-perfect}
\externaldocument[spaces-more-morphisms-]{spaces-more-morphisms}
\externaldocument[spaces-flat-]{spaces-flat}
\externaldocument[spaces-groupoids-]{spaces-groupoids}
\externaldocument[spaces-more-groupoids-]{spaces-more-groupoids}
\externaldocument[bootstrap-]{bootstrap}
\externaldocument[spaces-pushouts-]{spaces-pushouts}
\externaldocument[spaces-chow-]{spaces-chow}
\externaldocument[groupoids-quotients-]{groupoids-quotients}
\externaldocument[spaces-more-cohomology-]{spaces-more-cohomology}
\externaldocument[spaces-simplicial-]{spaces-simplicial}
\externaldocument[spaces-duality-]{spaces-duality}
\externaldocument[formal-spaces-]{formal-spaces}
\externaldocument[restricted-]{restricted}
\externaldocument[spaces-resolve-]{spaces-resolve}
\externaldocument[formal-defos-]{formal-defos}
\externaldocument[defos-]{defos}
\externaldocument[cotangent-]{cotangent}
\externaldocument[examples-defos-]{examples-defos}
\externaldocument[algebraic-]{algebraic}
\externaldocument[examples-stacks-]{examples-stacks}
\externaldocument[stacks-sheaves-]{stacks-sheaves}
\externaldocument[criteria-]{criteria}
\externaldocument[artin-]{artin}
\externaldocument[quot-]{quot}
\externaldocument[stacks-properties-]{stacks-properties}
\externaldocument[stacks-morphisms-]{stacks-morphisms}
\externaldocument[stacks-limits-]{stacks-limits}
\externaldocument[stacks-cohomology-]{stacks-cohomology}
\externaldocument[stacks-perfect-]{stacks-perfect}
\externaldocument[stacks-introduction-]{stacks-introduction}
\externaldocument[stacks-more-morphisms-]{stacks-more-morphisms}
\externaldocument[stacks-geometry-]{stacks-geometry}
\externaldocument[moduli-]{moduli}
\externaldocument[moduli-curves-]{moduli-curves}
\externaldocument[examples-]{examples}
\externaldocument[exercises-]{exercises}
\externaldocument[guide-]{guide}
\externaldocument[desirables-]{desirables}
\externaldocument[coding-]{coding}
\externaldocument[obsolete-]{obsolete}
\externaldocument[fdl-]{fdl}
\externaldocument[index-]{index}

% Theorem environments.
%
\theoremstyle{plain}
\newtheorem{theorem}[subsection]{Theorem}
\newtheorem{proposition}[subsection]{Proposition}
\newtheorem{lemma}[subsection]{Lemma}

\theoremstyle{definition}
\newtheorem{definition}[subsection]{Definition}
\newtheorem{example}[subsection]{Example}
\newtheorem{exercise}[subsection]{Exercise}
\newtheorem{situation}[subsection]{Situation}

\theoremstyle{remark}
\newtheorem{remark}[subsection]{Remark}
\newtheorem{remarks}[subsection]{Remarks}

\numberwithin{equation}{subsection}

% Macros
%
\def\lim{\mathop{\mathrm{lim}}\nolimits}
\def\colim{\mathop{\mathrm{colim}}\nolimits}
\def\Spec{\mathop{\mathrm{Spec}}}
\def\Hom{\mathop{\mathrm{Hom}}\nolimits}
\def\Ext{\mathop{\mathrm{Ext}}\nolimits}
\def\SheafHom{\mathop{\mathcal{H}\!\mathit{om}}\nolimits}
\def\SheafExt{\mathop{\mathcal{E}\!\mathit{xt}}\nolimits}
\def\Sch{\mathit{Sch}}
\def\Mor{\mathop{\mathrm{Mor}}\nolimits}
\def\Ob{\mathop{\mathrm{Ob}}\nolimits}
\def\Sh{\mathop{\mathit{Sh}}\nolimits}
\def\NL{\mathop{N\!L}\nolimits}
\def\CH{\mathop{\mathrm{CH}}\nolimits}
\def\proetale{{pro\text{-}\acute{e}tale}}
\def\etale{{\acute{e}tale}}
\def\QCoh{\mathit{QCoh}}
\def\Ker{\mathop{\mathrm{Ker}}}
\def\Im{\mathop{\mathrm{Im}}}
\def\Coker{\mathop{\mathrm{Coker}}}
\def\Coim{\mathop{\mathrm{Coim}}}

% Boxtimes
%
\DeclareMathSymbol{\boxtimes}{\mathbin}{AMSa}{"02}

%
% Macros for moduli stacks/spaces
%
\def\QCohstack{\mathcal{QC}\!\mathit{oh}}
\def\Cohstack{\mathcal{C}\!\mathit{oh}}
\def\Spacesstack{\mathcal{S}\!\mathit{paces}}
\def\Quotfunctor{\mathrm{Quot}}
\def\Hilbfunctor{\mathrm{Hilb}}
\def\Curvesstack{\mathcal{C}\!\mathit{urves}}
\def\Polarizedstack{\mathcal{P}\!\mathit{olarized}}
\def\Complexesstack{\mathcal{C}\!\mathit{omplexes}}
% \Pic is the operator that assigns to X its picard group, usage \Pic(X)
% \Picardstack_{X/B} denotes the Picard stack of X over B
% \Picardfunctor_{X/B} denotes the Picard functor of X over B
\def\Pic{\mathop{\mathrm{Pic}}\nolimits}
\def\Picardstack{\mathcal{P}\!\mathit{ic}}
\def\Picardfunctor{\mathrm{Pic}}
\def\Deformationcategory{\mathcal{D}\!\mathit{ef}}
\usepackage{cleveref}
\crefname{originaltheorem}{original theorem}{original theorems}
\Crefname{originaltheorem}{Original theorem}{Original theorems}
\makeatletter
\newlabel{theorem-bondal-van-den-Bergh}{{09IS}{source}{}{sourcepointer-1}{}}
\newlabel{theorem-bondal-van-den-Bergh@cref}{{[originaltheorem][0][]09IS}{[1][0][]source}{}{}{}}
\makeatother
\crefname{originalsection}{original section}{original sections}
\Crefname{originalsection}{Original section}{Original sections}
\makeatletter
\newlabel{dga-section-variant-base-change}{{09LU}{source}{}{sourcepointer-2}{}}
\newlabel{dga-section-variant-base-change@cref}{{[originalsection][0][]09LU}{[1][0][]source}{}{}{}}
\makeatother
\crefname{originallemma}{original lemma}{original lemmas}
\Crefname{originallemma}{Original lemma}{Original lemmas}
\makeatletter
\newlabel{lemma-ext-from-perfect-into-bounded-QCoh}{{09M4}{source}{}{sourcepointer-3}{}}
\newlabel{lemma-ext-from-perfect-into-bounded-QCoh@cref}{{[originallemma][0][]09M4}{[1][0][]source}{}{}{}}
\makeatother
\crefname{originallemma}{original lemma}{original lemmas}
\Crefname{originallemma}{Original lemma}{Original lemmas}
\makeatletter
\newlabel{dga-lemma-tensor-with-complex-derived}{{09LX}{source}{}{sourcepointer-4}{}}
\newlabel{dga-lemma-tensor-with-complex-derived@cref}{{[originallemma][0][]09LX}{[1][0][]source}{}{}{}}
\makeatother
\crefname{originalproposition}{original proposition}{original propositions}
\Crefname{originalproposition}{Original proposition}{Original propositions}
\makeatletter
\newlabel{proposition-compact-is-perfect}{{09M1}{source}{}{sourcepointer-5}{}}
\newlabel{proposition-compact-is-perfect@cref}{{[originalproposition][0][]09M1}{[1][0][]source}{}{}{}}
\makeatother
\crefname{originallemma}{original lemma}{original lemmas}
\Crefname{originallemma}{Original lemma}{Original lemmas}
\makeatletter
\newlabel{dga-lemma-fully-faithful-in-compact-case}{{09LZ}{source}{}{sourcepointer-6}{}}
\newlabel{dga-lemma-fully-faithful-in-compact-case@cref}{{[originallemma][0][]09LZ}{[1][0][]source}{}{}{}}
\makeatother
\crefname{originallemma}{original lemma}{original lemmas}
\Crefname{originallemma}{Original lemma}{Original lemmas}
\makeatletter
\newlabel{dga-lemma-tensor-with-complex-hom-adjoint}{{09LY}{source}{}{sourcepointer-7}{}}
\newlabel{dga-lemma-tensor-with-complex-hom-adjoint@cref}{{[originallemma][0][]09LY}{[1][0][]source}{}{}{}}
\makeatother
\crefname{originallemma}{original lemma}{original lemmas}
\Crefname{originallemma}{Original lemma}{Original lemmas}
\makeatletter
\newlabel{lemma-tensor-with-QCoh-complex}{{09M3}{source}{}{sourcepointer-8}{}}
\newlabel{lemma-tensor-with-QCoh-complex@cref}{{[originallemma][0][]09M3}{[1][0][]source}{}{}{}}
\makeatother
\crefname{originallemma}{original lemma}{original lemmas}
\Crefname{originallemma}{Original lemma}{Original lemmas}
\makeatletter
\newlabel{derived-lemma-fully-faithful-adjoint-kernel-zero}{{09J1}{source}{}{sourcepointer-9}{}}
\newlabel{derived-lemma-fully-faithful-adjoint-kernel-zero@cref}{{[originallemma][0][]09J1}{[1][0][]source}{}{}{}}
\makeatother
\title{The derived category of quasi-coherent modules on a quasi-compact quasi-separated scheme is a derived category of a differential graded algebra}
\author{The Stacks Project Authors}
\date{stacks-project:perfect.tex:local-snapshot-20261012; GNU FDL 1.2 or later}
\begin{document}
\maketitle
\setcounter{subsection}{1}
\noindent\textit{Editorial note.} All delivered mathematical text is the author's own, copied from the pinned Stacks Project chapter file perfect.tex; no statement or proof step was rewritten, reordered or removed. The counted span is the author environment and its own proof environment. References inside the span are handed to the assembler as navigation stubs that carry the canonical tag of the cited tag only; whether each of them is load-bearing is recorded on the stub itself.\par\medskip
\noindent\textit{Allowed standard prerequisites (editorial).}\begin{itemize}
\item A quasi-compact quasi-separated scheme admits a perfect K-injective complex that generates the derived category of quasi-coherent modules, by the Bondal-van den Bergh theorem and the existence of K-injective resolutions. invoked at proof lines 4690-4693. Provides the generator K\textasciicircum{}bullet out of which the dg algebra is built.
\item For a K-injective complex the cohomology of the dg endomorphism algebra computes the Ext groups of the corresponding object. invoked at proof lines 4701-4706. Identifies H\textasciicircum{}n(E) with Ext\textasciicircum{}n(K\textasciicircum{}bullet, K\textasciicircum{}bullet), whose finiteness follows from the vanishing lemma for Ext out of perfect objects into bounded quasi-coherent complexes.
\item The dg tensor functor - tensor\textasciicircum{}L\_E K\textasciicircum{}bullet and its right adjoint R Hom(K\textasciicircum{}bullet, -), with the fully faithfulness criterion in the compact case and the equivalence criterion for a fully faithful functor whose adjoint has vanishing kernel. invoked at proof lines 4715-4727 and proof lines 4737-4739. Assembles the equivalence from the adjunction.
\end{itemize}\medskip
\ifdefined\hypertarget\hypertarget{sourcepointer-1}{}\fi
\noindent\textit{Source reference.} Original theorem 09IS: sources/stacks-perfect/pinned-original-source.tex, source lines 4111--4111. Cited at proof lines 4692-4693 for the existence of K\textasciicircum{}bullet; the construction of (E, d) from it is inside the counted span.\par\medskip
\ifdefined\hypertarget\hypertarget{sourcepointer-2}{}\fi
\noindent\textit{Source reference.} Original section 09LU: sources/stacks-chow/tags.txt, source lines 12157--12157. Cited at proof line 4700 as the reference section for the dg category conventions; the definition of (E, d) is inside the counted span.\par\medskip
\ifdefined\hypertarget\hypertarget{sourcepointer-3}{}\fi
\noindent\textit{Source reference.} Original lemma 09M4: sources/stacks-perfect/pinned-original-source.tex, source lines 3465--3465. Cited at proof lines 4707-4708 for finiteness of the nonzero cohomology groups of E; the Ext identification is inside the counted span.\par\medskip
\ifdefined\hypertarget\hypertarget{sourcepointer-4}{}\fi
\noindent\textit{Source reference.} Original lemma 09LX: sources/stacks-chow/tags.txt, source lines 12160--12160. Cited at proof lines 4715-4716 to define the tensor functor; the fully faithfulness argument is inside the counted span.\par\medskip
\ifdefined\hypertarget\hypertarget{sourcepointer-5}{}\fi
\noindent\textit{Source reference.} Original proposition 09M1: sources/stacks-perfect/pinned-original-source.tex, source lines 4464--4464. Cited at proof lines 4717-4718 for the compactness of K\textasciicircum{}bullet; the fully faithfulness conclusion is inside the counted span.\par\medskip
\ifdefined\hypertarget\hypertarget{sourcepointer-6}{}\fi
\noindent\textit{Source reference.} Original lemma 09LZ: sources/stacks-chow/tags.txt, source lines 12162--12162. Cited at proof lines 4721-4722 for full faithfulness; the compactness input and the Ext identification are inside the counted span.\par\medskip
\ifdefined\hypertarget\hypertarget{sourcepointer-7}{}\fi
\noindent\textit{Source reference.} Original lemma 09LY: sources/stacks-chow/tags.txt, source lines 12161--12161. Cited at proof lines 4726-4727 for the right adjoint; the quasi-inverse and kernel arguments are inside the counted span.\par\medskip
\ifdefined\hypertarget\hypertarget{sourcepointer-8}{}\fi
\noindent\textit{Source reference.} Original lemma 09M3: sources/stacks-perfect/pinned-original-source.tex, source lines 4650--4650. Cited at proof line 4730 to land the tensor functor in D\_QCoh; the kernel vanishing by generation is inside the counted span.\par\medskip
\ifdefined\hypertarget\hypertarget{sourcepointer-9}{}\fi
\noindent\textit{Source reference.} Original lemma 09J1: sources/stacks-chow/tags.txt, source lines 12059--12059. Cited at proof lines 4738-4739 to conclude the equivalence; the kernel vanishing by generation is inside the counted span.\par\medskip
\begin{theorem}
\label{theorem-DQCoh-is-Ddga}
Let $X$ be a quasi-compact and quasi-separated scheme.
Then there exist a differential graded algebra $(E, \text{d})$
with only a finite number of nonzero cohomology groups $H^i(E)$
such that $D_\QCoh(\mathcal{O}_X)$ is equivalent
to $D(E, \text{d})$.
\end{theorem}

\begin{proof}
Let $K^\bullet$ be a K-injective complex of $\mathcal{O}$-modules which
is perfect and generates $D_\QCoh(\mathcal{O}_X)$. Such a
thing exists by Theorem \ref{theorem-bondal-van-den-Bergh}
and the existence of K-injective resolutions. We will show the
theorem holds with
$$
(E, \text{d}) = \Hom_{\text{Comp}^{dg}(\mathcal{O}_X)}(K^\bullet, K^\bullet)
$$
where $\text{Comp}^{dg}(\mathcal{O}_X)$ is the differential graded category
of complexes of $\mathcal{O}$-modules. Please see
Differential Graded Algebra, Section \ref{dga-section-variant-base-change}.
Since $K^\bullet$ is K-injective we
have
\begin{equation}
\label{equation-E-is-OK}
H^n(E) = \Ext^n_{D(\mathcal{O}_X)}(K^\bullet, K^\bullet)
\end{equation}
for all $n \in \mathbf{Z}$. Only a finite number of these Exts
are nonzero by Lemma \ref{lemma-ext-from-perfect-into-bounded-QCoh}.
Consider the functor
$$
- \otimes_E^\mathbf{L} K^\bullet :
D(E, \text{d}) \longrightarrow D(\mathcal{O}_X)
$$
of
Differential Graded Algebra, Lemma
\ref{dga-lemma-tensor-with-complex-derived}.
Since $K^\bullet$ is perfect, it defines a compact object of
$D(\mathcal{O}_X)$, see Proposition \ref{proposition-compact-is-perfect}.
Combined with (\ref{equation-E-is-OK}) the functor above is fully
faithful as follows from
Differential Graded Algebra, Lemmas
\ref{dga-lemma-fully-faithful-in-compact-case}. It has a right adjoint
$$
R\Hom(K^\bullet, - ) : D(\mathcal{O}_X) \longrightarrow D(E, \text{d})
$$
by Differential Graded Algebra, Lemmas
\ref{dga-lemma-tensor-with-complex-hom-adjoint}
which is a left quasi-inverse functor by generalities on adjoint
functors. On the other hand, it follows from
Lemma \ref{lemma-tensor-with-QCoh-complex} that we obtain
$$
- \otimes_E^\mathbf{L} K^\bullet :
D(E, \text{d}) \longrightarrow D_\QCoh(\mathcal{O}_X)
$$
and by our choice of $K^\bullet$ as a generator of
$D_\QCoh(\mathcal{O}_X)$ the kernel of the adjoint
restricted to $D_\QCoh(\mathcal{O}_X)$ is zero.
A formal argument shows that we obtain the desired equivalence, see
Derived Categories, Lemma
\ref{derived-lemma-fully-faithful-adjoint-kernel-zero}.
\end{proof}

\end{document}
